\documentclass[
    prd, superscriptaddress, nofootinbib, amsmath, amssymb,
    aps, floatfix, preprintnumbers
]{revtex4-2}

\usepackage[utf8]{inputenc}
\usepackage{setspace}
\usepackage[dvips]{graphicx}
\usepackage[dvipsnames]{xcolor}
\definecolor{CiteBlue}{RGB}{45,52,151}
\usepackage[
    colorlinks=true,
    linkcolor=CiteBlue,
    urlcolor=CiteBlue,
    citecolor=CiteBlue
]{hyperref}

\usepackage[capitalise]{cleveref}
\usepackage{siunitx}

\usepackage{booktabs}

\newcommand{\refcite}[1]{Ref.~\cite{#1}}
\newcommand{\refscite}[1]{Refs.~\cite{#1}}

\usepackage{bm}
\newcommand{\bb}[1]{\bm{\mathrm{#1}}}
\newcommand{\du}{\mathrm{d}}
\newcommand{\dd}{\,\du}

\usepackage{verbatim}

\begin{document}
\title{Paper preparation pet peeves}
 
\author{Benjamin V. Lehmann}
\email{benvlehmann@gmail.com}
\affiliation{Center for Theoretical Physics---a Leinweber Institute, Massachusetts Institute of Technology, Cambridge, MA 02139, USA}

\date\today


\begin{abstract}\ignorespaces{}
\end{abstract}

\maketitle

\section{Labels and references}
\label{sec:labels-and-references}
\begin{figure}\centering
    \caption{This is a figure caption. \textbf{Do not} use \texttt{[h]}, \texttt{[t]}, or any other modifiers.}
    \label{fig:my-figure}
\end{figure}
\begin{enumerate}
    \item \textbf{Do} begin labels with \texttt{sec:}, \texttt{eq:}, \texttt{fig:}, \texttt{tab:}, etc., as appropriate. \textbf{Never} put spaces in labels.
    \item \textbf{Do not} use \verb|\ref| or \verb|\eqref|. Use \verb|\cref|. 
        \begin{itemize}
            \item Example: \cref{fig:my-figure}.
            \item At the beginning of a sentence, use \verb|\Cref|, which produces e.g.\ ``\Cref{fig:my-figure}.''
        \end{itemize}
    \item \textbf{Do not} use separate \verb|\cref| statements to reference multiple equations at once, e.g., ``\cref{eq:multi-line-equation} and \cref{eq:piecewise-expression}.'' Instead, use a single statement, as in ``\cref{eq:multi-line-equation,eq:piecewise-expression}.''
    \item \textbf{If and only if} a citation is part of a sentence, describe it as ``\refcite{Essig:2011nj},'' not ``\cite{Essig:2011nj}.'' Pluralize as ``\refscite{Essig:2011nj,Chandrasekhar:1942}.''
\end{enumerate}

\section{Units}
\begin{enumerate}
    \item \textbf{Do not} use \verb|\mathrm| for units. Use \verb|\qty| from the \texttt{siunitx} package. Examples: \qty{10}{\mega\electronvolt/\centi\meter^2}, \qty{2.2e3}{M_\odot}.
    \item To use the same syntax for dimensionless numbers, use \verb|\num|: \num{4e7}, \num{e3}, \num{427767865863}.
\end{enumerate}

\section{Equations}
\begin{enumerate}
    \item \textbf{Do not} use \verb|$$...$$| or \verb|\[...\]| for display equations. Use the \texttt{equation} environment.
    \item \textbf{Do not} use \texttt{align} for single equations that need to be split over multiple lines. Use \texttt{multline}. Example:
    \begin{multline}
        \label{eq:multi-line-equation}
        v_p^{{\mathrm{ex}}}(\theta,\phi) = \sqrt{
            \frac1{\epsilon_\parallel} +
            \frac12\left(
                \frac{1}{\epsilon_\perp} - \frac{1}{\epsilon_\parallel}
            \right)\Bigl[
            (\cos\phi\sin\theta)\sqrt{u_-} - (\cos\theta) \sqrt{u_+}
        \Bigr]^2}
        \\
        \pm \frac{
            \epsilon_\perp - \epsilon_\parallel
            + \epsilon_\perp \epsilon_\parallel v^2\cos2\psi
        }{
            \epsilon_\perp \epsilon_\parallel v^2\sqrt{
                \left[
                    \frac{1}{v^2}\left(
                            1/\epsilon_\parallel
                            - 1/\epsilon_\perp
                    \right) + \cos2\psi
                \right]^2 + \sin^2 2\psi
            }
        }
        .
    \end{multline}
    \item \textbf{Do not} use \texttt{eqnarray} or \texttt{array} for piecewise expressions. Use \texttt{cases} instead. Example:
    \begin{equation}
        \label{eq:piecewise-expression}
        P_{\mathrm{obs}}(\lambda) \approx \begin{cases}
            \displaystyle
            \lambda e^{-\lambda} P\bigl(\left|\delta E\right| > E_{\min}\bigr)
            & \lambda \ll 1, \\[3mm]\displaystyle
            \operatorname{erfc}\left(
                \frac{E_{\min}}{\sigma_E\sqrt{2\lambda}}\right)
            & \lambda \gg 1.
        \end{cases}
    \end{equation}
    \item \textbf{Do} use newlines to organize the code for long equations, as in the above examples. \textbf{Do not} exceed 80 characters per line in the code within the \texttt{equation} environment.
\end{enumerate}

\section{Tables and arrays}
\begin{table}
    \centering
    \begin{equation*}
        \begin{array}{lccc}
            \toprule
                & \textnormal{\textbf{Parameter}}
                    & \textnormal{\textbf{Loose}}
                    & \textnormal{\textbf{Tight}}
            \\\midrule
                \textnormal{\textbf{Cuts}}
                    & \tau_+~[\qty{}{\micro\second}]
                        & (0.3, 0.4) & (0.325, 0.375) \\
                    & \tau_-~[\qty{}{\micro\second}]
                        & (1.6, 2.0) & (1.75, 1.85) \\
                    & \chi^2_\text{red}
                        & (0.9, 1.1)& (0.95, 1.05) \\
                    & t_0~[\qty{}{\micro\second}]
                        & (-0.5, 0.5) & (-0.5, 0.5)
            \\\midrule
                \textnormal{\textbf{Results}}
                    & \textnormal{Event rate}~[\qty{}{\hertz}]
                        & \num{7.2e-5} & \num{7.4e-6} \\
                    & \textnormal{Survival}~[\qty{}{\percent}]
                        & 0.107 & 0.011 \\
            \bottomrule
        \end{array}
    \end{equation*}
    \caption{This is a table consisting mainly of mathematical expressions.}
    \label{tab:math}
\end{table}
\begin{table}
    \centering
    \setlength{\tabcolsep}{6pt}
    \begin{tabular}{lcccccccl}
        \toprule
        System
            & $m_1$ [\qty{}{M_\odot}]
            & $m_2$ [\qty{}{M_\odot}]
            & $a_{\mathrm{proj}}$ [\qty{}{AU}]
            & $e$
            & $\theta_i$ [\qty{}{\degree}]
            & $\dot P_{12}$
            & $\Delta P_{12}/m_3$ [\qty{}{yr/M_\odot}]
            & Ref. \\
        \midrule
        J1713-0747
            & 1.33 & 0.29 & 0.0648 & \num{7.49e-5}
            & 71.69 & \num{3.4e-13} & 0.24014
            & \cite{Zhu:2018etc} \\
        J1740-3052
            & 1.4 & 19.5 & 1.52 & 0.5789
            & 53 & \num{3e-9} & 0.0804596
            & \cite{Madsen:2012rs} \\
        J1903+0327
            & 1.667 & 1.029 & 0.211 & 0.4367
            & 77.47 & \num{-3.3e-11} & 0.229562
            & \cite{Freire:2010tf} \\
        J2016+1948
            & ${\sim} 1$ & 0.29 & 0.302 & \num{1.48e-3}
            & 58.58 & -- & 1.41564
            & \cite{Gonzalez:2011kt} \\
        B1259-63
            & ${\sim} 2$ & 19.8 & \num{2.60} & 0.8698
            & $154 \pm 3$  & \num{1.4e-8} &  0.304914
            & \cite{Miller-Jones:2018waj, Shannon:2013dpa} \\
        \bottomrule
    \end{tabular}
    \caption{This is a table containing significant nonmathematical content.}
    \label{tab:nonmath}
\end{table}

\begin{enumerate}
    \item When a table consists largely of numerical or mathematical expressions, use the \verb|array| environment within an appropriate math environment. See e.g.\ \cref{tab:math}.
    \item When a table has significant nonmathematical content that would be inappropriate for a mathmode environment, use \verb|tabular|, as in \cref{tab:nonmath}.
    \item Avoid vertical rules in tables. For horizontal rules, use \verb|\toprule|, \verb|\midrule|, and \verb|\bottomrule| from the \texttt{booktabs} package, as in \cref{tab:math,tab:nonmath}.
\end{enumerate}

\section{Text}
\begin{enumerate}
    \item \textbf{Do not} use \verb|{\bf x}|, \verb|{\it x}|, \verb|{\rm x}|, or \verb|{\color{x} y}|. Instead, use \verb|\textbf{x}|, \verb|\textit{x}|, \verb|\mathrm{x}|, and \verb|\textcolor{x}{y}|.
    \item \textbf{Do not} use \verb|\mathrm| for text with spaces or dashes in math mode. Instead, use \verb|\textnormal|: $x\textnormal{--}y_{\textnormal{(a b)}}$.
    \item \textbf{Do not} use \verb|%| over more than a few lines. Use the \texttt{comment} environment provided by the \texttt{verbatim} package.
\end{enumerate}

\section{Dashes}
\begin{enumerate}
    \item Use \verb|-| (hyphen) for hyphenation, as in the word ``non-negotiable.'' Do not surround it with spaces.
    \item Use \verb|--| (en dash) to indicate ranges, as in ``1--10.'' Do not surround it with spaces.\footnote{My affiliation breaks this rule, which breaks my heart.}
    \item Use \verb|---| (em dash) as a punctuation mark---like this. Do not surround it with spaces.
\end{enumerate}

\section{Spacing}
\begin{enumerate}
    \item \textbf{Do not} use an ordinary space where a line break would be problematic (e.g., between a phrase and an associated \verb|\cite| invocation). Instead, use \verb|~|.
    \item \textbf{Do not} use \verb|~| in math mode. Instead, use \verb|\,|, \verb|\;|, \verb|\quad|, \verb|\qquad|, or \verb|\hspace|, as appropriate.
    \item When a period is \textbf{not} used to end a sentence, and immediately precedes a space, either use an explicit interword space, ``\verb|\ |'', or write ``\verb|\@.|'' instead of ``\verb|.|''. See e.g.\ the usage of ``e.g.'' throughout this document. The same applies for semicolons that precede a space but do not end a clause.
\end{enumerate}

\section{Miscellaneous}
\begin{enumerate}
\item \textbf{Do not} use custom macros to begin and end environments, particularly equation environments like \texttt{equation} and \texttt{align}. This makes debugging more complicated, and also interferes with syntax highlighting and linting. If a line begins an environment of any kind, it should say \verb|\begin|. If you need more functionality than the environment provides, define a new environment.
\item \textbf{Do not} start an environment in the middle of a line if avoidable. Always allow \verb|\begin{...}| and \verb|\end{...}| to sit on their own lines. \text{Do not} add blank lines before or after the environment unless you mean to make a paragraph break.
\item \textbf{Do not} use binary operator spacing for \verb|\sim| when using it as a unary operator. Instead, use \verb|\mathop{\sim}| or simply \verb|{\sim}|. For example, notice the difference between $\sim\qty{e3}{\centi\meter}$, $\mathop{\sim}\qty{e3}{\centi\meter}$, and ${\sim}\qty{e3}{\centi\meter}$, which are especially pronounced when a line has extra space before a newline.
\end{enumerate}

\bibliography{style}

\end{document}
